\chapter{Control Flow and Dominance}
\label{ch:pebble-dominance}

An ID tells us where a value is stored in the compiler, not whether its definition executes before a use.
Control flow supplies the missing condition.
PebbleIR derives its graph from terminators rather than maintaining a second editable graph that could disagree with the instructions.

\section{One source of truth for edges}
\label{sec:pebble-edges}
A return has no successors, a jump has one, and a branch has two distinct successors.
Invert these edges to obtain predecessor sets.
The verifier stores one bit per predecessor in an unsigned word; the eight-block limit makes the shifts and masks fit even on an implementation with 16-bit unsigned integers.
The graph contains only declared blocks and validated targets before any traversal begins.

Figure~\ref{fig:pebble-diamond} is the graph of the executable diamond fixture.
The left arm returns an incremented input to the join; the right arm returns a doubled input.
The join's phi selects the arm that actually ran.

\begin{figure}[htbp]
\centering
\begin{tikzpicture}[node distance=13mm and 15mm]
\node[basic block] (e) {0: test x < 10};
\node[basic block,below left=of e] (l) {1: left = x + 1};
\node[basic block,below right=of e] (r) {2: right = x * 2};
\node[basic block,below right=of l] (j) {3: phi; return};
\draw[cfg edge] (e)--node[above left]{true}(l);
\draw[cfg edge] (e)--node[above right]{false}(r);
\draw[cfg edge] (l)--(j); \draw[cfg edge] (r)--(j);
\end{tikzpicture}
\caption{Neither arm dominates the join, but entry dominates every block.}
\label{fig:pebble-diamond}
\end{figure}

Reachability begins at block zero and follows successor relationships to closure.
An unreachable block can have predecessors; an unreachable cycle is the simplest counterexample to an in-degree test.
Our verifier rejects all unreachable declared blocks.
A future pass that folds a branch must therefore remove newly unreachable blocks and repair phis before it publishes a verified function.
The current optimizer avoids that extra obligation by retaining the original CFG.

\section{Compute the property the client needs}
\label{sec:pebble-dom-equations}
A block $D$ \term{dominates} a block $B$ when every entry-to-$B$ path contains $D$.
The equations are $\dom(0)=\{0\}$ and, for other blocks,
\[
  \dom(B)=\{B\}\cup\bigcap_{P\in\Pred(B)}\dom(P).
\]
Algorithm~\ref{alg:pebble-dom} computes their greatest fixed point by beginning with an overestimate.
The entry's fixed boundary is not recomputed from an empty predecessor intersection.

\begin{algorithm}[htbp]
\caption{Dominator sets for a verified, entry-reachable graph.}
\label{alg:pebble-dom}
\begin{pseudocode}
\Require reachable blocks $V$, predecessor sets, predecessor-free entry 0
\Ensure dominator set for every block
\State $D[0]\gets\{0\}$; $D[B]\gets V$ for $B\ne0$
\Repeat
\State $changed\gets false$
\ForAll{$B\ne0$}
\State $new\gets\{B\}\cup\bigcap_{P\in\Pred(B)}D[P]$
\If{$new\ne D[B]$}
\State $D[B]\gets new$; $changed\gets true$
\EndIf
\EndFor
\Until{not $changed$}
\end{pseudocode}
\end{algorithm}

Each strict update removes at least one candidate dominator, so the finite process terminates.
At termination the equations hold; the standard path characterization identifies this greatest solution with dominance on the reachable graph.
Starting every set empty would select the wrong solution in the presence of cycles.
Do not substitute a different initialization merely because the iteration still terminates.

\section{Work the diamond and the loop}
\label{sec:pebble-dom-worked}
For the diamond, the predecessor sets are empty, $\{0\}$, $\{0\}$, and $\{1,2\}$.
The final dominators are respectively $\{0\}$, $\{0,1\}$, $\{0,2\}$, and $\{0,3\}$.
At the join, intersecting $\{0,1\}$ and $\{0,2\}$ retains only entry; adding the join itself yields $\{0,3\}$.
Thus neither arm's result can be used as an ordinary join operand without a phi.

For the loop graph $0\to1$, $1\to2$, $2\to1$, and $1\to3$, the final sets are $\{0\}$, $\{0,1\}$, $\{0,1,2\}$, and $\{0,1,3\}$.
The header dominates the body and exit.
A header phi may consequently use its own old value as an incoming operand from the body, provided the edge-use rule is satisfied.
This property is used by the parallel-swap fixture.

\section{Representation cost and invalidation}
\label{sec:pebble-dom-cost}
One sweep scans predecessor candidates for every nonentry block, costing $O(B^2)$ with one-word masks.
For $R$ sweeps the cost is $O(RB^2)$.
These small bounds describe the implementation, not an asymptotically competitive dominator algorithm for huge functions.
A larger graph should use adjacency lists and a suitable dominator algorithm while preserving the same query meaning.

Changing any terminator can change reachability, predecessors, and dominance.
Even a value-only edit changes the function revision in the conservative cache policy.
That policy sometimes recomputes unchanged graph facts, but it is easier to audit than a preservation declaration that quietly misses one edge mutation.

\section{Exercises}
\label{sec:pebble-nine-exercises}
\begin{enumerate}
\item \label{ex:pebble-shortcut} Add edge $0\to3$ to the diamond as a graph exercise. Recompute $\dom(3)$ and explain why a previously two-input join phi needs repair.
\item \label{ex:pebble-unreachable-cycle} Give predecessor sets for an unreachable two-block cycle and explain why neither block has in-degree zero.
\end{enumerate}

\section{Further reading}
\label{sec:pebble-nine-reading}
Cytron and colleagues connect dominance to SSA placement and control dependence~\cite{cytron-1991-ssa}.
Kildall supplies the broader dataflow perspective behind iterative equation solving~\cite{kildall-1973-dataflow}.
